
\lesson{04/11}
\section{Gibbs' Entropy functional}
[buh]

$$S[p(\nu)] = -k\sum_{\nu}^{}P(\nu)\ln P(\nu)$$
So the functional is extensive
$$S[P_{AB}(\nu_A,\nu_B)] = S_A+S_B$$
\subsection{Microcanonical ensemble}
The only constraint is
$$\sum_{\nu}^{}P(\nu) = 1,$$
for the probabilistic interpretation. As always, the maximization of entropy
will be constrained with \textit{Lagrange multipliers}.

Define
$$\tilde S [P]=\underbrace{-k\sum_{\nu}^{}P(\nu)\ln P(\nu)}_{=S[P]}+\alpha_0
\left(\sum_{\nu}^{}P(\nu) -1\right)$$
and extremize it with arbitrary variations:
$$\delta\tilde S[P] = 0 = -k\sum_{\nu}^{}\ln P(\nu)\delta P(\nu) - k
\sum_{\nu}^{}P(\nu)\frac{\delta P(\nu)}{P(\nu)} + \alpha_0 \sum_{\nu}^{}\delta
P(\nu) = -k\sum_{\nu}^{}\left(-k\ln P(\nu) + (\alpha_0-k)\right)\delta P(\nu)
\footnote{occhio ai segni}$$
Since the variation of $P$ is arbitrary, this is null if and only if
$$k\ln P(\nu) = (\alpha_0-k) \iff P(\nu) = \exp\left[\frac{\alpha_0-k}{k}
\right]:$$
since this is constant with respect to $\nu$, the \textit{microstates are
equiprobable} (as found before); imposing the constraint
$$\sum_{\nu}^{}P(\nu) = P\sum_{\nu}^{}1 = \Omega e^{\alpha_0/k - 1}=1
\iff \frac{1}{\Omega} = e^{\alpha_0/k-1}$$
\subsection{Canonical ensemble}
Here the constraints are
$$\sum_{\nu}^{}P(\nu) = 1,\; \langle E\rangle_P = \sum_{\nu}^{}E(\nu)P(\nu) = E$$
Now the functional to maximise is
$$\tilde S[P] = -k\sum_{\nu}^{}P(\nu)\ln P(\nu) +\alpha_0\left(\sum_{\nu}^{}P(\nu)
-1\right) + \alpha_E\left(\sum_{\nu}^{}P(\nu)E(\nu)-E\right)$$
so
$$0=\delta\tilde S[P] = -k\sum_{\nu}^{}\ln P(\nu)\delta P(\nu) - k\sum_{\nu}^{}
P(\nu)\frac{\delta P(\nu)}{P(\nu)} + \alpha_0\sum_{\nu}^{}\delta P(\nu) + 
\alpha_E\sum_{\nu}^{}E(\nu)\delta P(\nu)
$$
writing it as a single sum\footnote{OCCHIO al segno di $\alpha_E$}
$$0 = \sum_{\nu}^{}\left[-k\ln P(\nu) +(\alpha_0-k)+\alpha_E E(\nu)\right]\delta
P(\nu),\,\forall \delta P(\nu)$$
So finally
$$\ln P(\nu) = \frac{\alpha_0}{k}-1 + \frac{\alpha_E}{k}E(\nu)\iff
P(\nu) = \exp\left[\frac{\alpha_0}{k}-1 + \frac{\alpha_E}{k}E(\nu)\right].
$$
Now the normalization constraint
$$1=\sum_{\nu}^{}P(\nu)=e^{\frac{\alpha_0}{k}-1}\sum_{\nu}^{}
e^{\flatfrac{\alpha_E}{k}E(\nu)}\iff Z := e^{\flatfrac{\alpha_0}{k}-1} = 
\sum_{\nu}^{}e^{\alpha_E/kE(\nu)}
$$
$Z$ will be called the \textit{partition function}.

With this notation, the entropy functional becomes
$$S[P(\nu)] = \sum_{\nu}^{}P(\nu)\left(\frac{\alpha_E}{k}E(\nu)-k\ln Z\right) = 
\alpha_E \langle E \rangle_P + k\ln Z
$$

The second constraint imposes
$$E=\sum_{\nu}^{}E(\nu)P(\nu)$$
since $E$ is the TD energy, the thermal equation of state must hold
$$\frac{1}{T} = \pdv[]{S}{\langle E\rangle} = -\alpha_E$$

So the microstate probability is
$$P(\nu) = \frac{1}{Z}e^{-\frac{1}{kT}E(\nu)} = \frac{1}{Z}e^{-\beta E(\nu)}$$
where
$$Z = \sum_{\nu}^{}e^{-\beta E(\nu)}.$$
Where are the TD potentials? Starting from the entropy
$$S = \frac{\langle E\rangle}{T}+\ln Z\iff\langle E\rangle -TS = -kT\ln Z$$
Since $E-TS:=F$, the Helmholtz potential, then
$$F = -kT\ln Z \iff Z = e^{-\beta F}$$
\subsection{Relation between microcanonical and canonical ensembles}
The key quantity in the canonical ensemble is the partition function
$$Z(T,V,N) = \sum_{\nu}^{}e^{-\beta E(\nu)}:$$
since the contributions of the various microstates are independent, the sum
can be reorganized: let $\Omega (E_e,V,N)$ be the \textit{number of microstates
of energy $E_e$}
$$Z = \sum_{e}^{}\sum_{\nu_e=1}^{\Omega(E_e)}e^{-\beta E_e} = \sum_{e}^{}
\Omega(E_e,V,N)e^{-\beta E_e}
$$
the $\Omega$ term is the key quantity of microcanonical ensembles and its weighted
sum over different energies is the partition function.

The continuous spectrum equivalent switches $\Omega$ with a density of states
$g(E)\dd E$ and the sum with an integral
$$Z(\beta) = \int_{0}^{\infty} \dd E g(E)e^{-\beta E}$$
The probability density $P(\nu)$ becomes
$$P(E)\dd E=\frac{1}{Z}e^{-\beta E} g(E)\dd E $$
\textbf{Note} $Z$ and $g(E)$ are \textit{not} independent and convey the same
information: they are related by a \textbf{Laplace transform}.

The Laplace anti-transform yields
$$g(E) = \frac{1}{2\pi i}\int_{\gamma}^{} \dd \beta Z(\beta) e^{\beta E}$$
where $\gamma$ is a contour in the complex space of $\beta$, parallel to the
imaginary axis and in the analyticity domain of $Z(\beta)$.

Since this transform relates the density of states via $g(E)$ with the
temperature via $Z(\frac{1}{kT})$, it is plausible that this expression,
in the TD limit, becomes the Legendre transform between $S$ and $T$:
since $S=k\ln\Omega$,
$$g(E)\to \eval{\Omega(E)}_{\Delta E=\dd E} = e^{S(E)/k} = g(E)\Delta E\implies
g(E)=\frac{e^{S(E)/k}}{\Delta E}$$
and
$$Z\to e^{-\beta F}$$
imply
$$e^{-\frac{F}{kT}} = \frac{1}{\Delta E}\int_{0}^{\infty} \dd E e^{\frac{S}{k}-
\frac{E}{kT}}$$
Since in the TD limit $S$ is large (it is extensive), this integral can be
computed using the \textbf{saddle-point technique}: the exponent
$$f(E) = \frac{1}{k}\left(S(E)-\frac{E}{T}\right) \to
0 = \pdv[]{f}{E} = \frac{1}{k}\left(\pdv[]{S}{E}-\frac{1}{T}\right)\iff
E\text{ is s.t. }\frac{1}{T} = \pdv[]{S}{E}$$
so the integrand is extremised around the TD energy $\langle E\rangle$!
The integral is therefore
$$e^{-\frac{F}{kT}} = e^{-f(\langle E\rangle)}\left[\frac{1}{\Delta E}
\sqrt{\frac{2\pi}{f''(\langle E\rangle)}} + ...\right]$$
Since $E$ is extensive $E\sim N$, so
$$f(E) \sim N \to f''\sim \frac{N}{N^2} = \frac{1}{N}$$
and so
$$\exp\left[-f(\langle E\rangle) - \frac{1}{2}\ln\left(\frac{\Delta E^2 2\pi}
{f''(\langle E\rangle)}\right)\right]\sim\exp\left[-f(\langle E\rangle)\right]$$
Finally
$$e^{-\frac{F}{kT}} = e^{-f(\langle E\rangle)} \iff -\frac{F}{T} = S-
\frac{\langle E\rangle}{T}\iff \boxed{F = \langle E\rangle-TS}$$
which is the TD Legendre transform.

Moreover, since the integrand is similar to $P(E)$, the TD energy is also
\textit{the most probable value of energy}.
\section{Canonical Treatment of non-interacting systems}
Suppose to have a system composed of $N$ distinguishable, non-interacting
components: each microstate of the system $\nu$ is a collection of $N$ microstates
$\left\{\nu_1,...\nu_N\right\}$ and its energy is simply
$$E_\nu = \sum_{i=1}^{N}E_{\nu_i}$$
and the partition function is
$$Z = \sum_{\nu}^{}e^{-\beta E_\nu} = \sum_{\nu_1}...\sum_{\nu_N}^{}
e^{-\beta E_{\nu_1}}...e^{-\beta E_{\nu_N}}:= Z_1...Z_N$$
where $Z_i$ is formally defined as the partition function for the $i$-th 
component.

Similarly, the expression of the Helmholtz energy is
$$[...\text{cancellata lavangna...} F=-kT\ln Z]$$

If the particles are indistinguishable, there is some ''over-counting``, so the
expression need be symmetrized and multiplied approximately by $\frac{1}{N!}$.
\subsection{The ideal classical gas (canonical ensemble)}
All the gas particles are identical, so $Z_i=Z_j$
$$Z = \frac{1}{N!}\prod_{i=1}^{N}Z_i = \frac{1}{N!}Z_1^{N}$$
Since $Z_1$ is for a single particle, it is easy to compute, since it is an
integral over its phase space
$$Z_1 = \int_{}^{} \frac{\dd ^3p\dd^3q}{h^3} e^{-\beta H(q,p)}=
\int_{}^{} \dd^3q\int_{}^{} \frac{\dd^3p}{h^3} e^{-\beta\frac{\mathbf{p}^2}{2m}}=
\frac{V}{h^3}(2\pi mkT)^{\frac{3}{2}}
$$
